Este trabalho teve como objetivo desenvolver um aplicativo de comunicação controlado pelo olhar que funcionasse em \textit{tablets} Android comuns, usando apenas a câmera frontal, sem \textit{hardware} adicional e sem conexão com a internet, como alternativa de baixo custo aos dispositivos comerciais. O objetivo foi alcançado: o aplicativo desenvolvido permite compor uma mensagem e falá-la usando somente o olhar, executa todo o processamento no próprio aparelho e foi usado de ponta a ponta em um Samsung Galaxy Tab S7, nas 15 sessões dos testes de uso descritos no \refCap{5_Resultados}.

Cada um dos objetivos específicos (Seção~\ref{sec:objetivos}) foi atendido. O rastreamento ocular combina a detecção dos cantos dos olhos pelo ML Kit, a localização da pupila por uma rede YOLO26n treinada para esse fim e uma calibração por usuário baseada em regressão polinomial, seguida de um filtro de Kalman (Seção~\ref{sec:met_eyetracking}). A interface de digitação converte a direção do olhar em comandos de navegação e seleção por meio de um tempo de permanência, e foi operada sem nenhum toque na tela. O preenchimento automático, baseado em trigramas com \textit{Stupid Backoff} e com tolerância a erros de digitação, sugere a palavra em andamento ou a próxima palavra a partir de um \textit{corpus} de frases de \gls{CAA} (Seção~\ref{sec:met_preenchimento}). A síntese de voz, feita pelo Kokoro executado localmente, converteu a mensagem em fala clara e fiel ao texto escrito em todas as sessões, sem envio de dados a servidores externos. Por fim, o módulo de métricas registrou todas as sessões de forma consistente, o que permitiu analisar a velocidade, o esforço e a contribuição das sugestões a partir dos próprios registros do aplicativo.

Os testes de uso levaram a três conclusões principais. A primeira é que o aplicativo pode ser aprendido rapidamente: nos três testes, o desempenho melhorou com a prática, de forma mais clara a partir da terceira iteração, e no T1 e no T2 a quinta iteração chegou ao número mínimo de comandos previsto pelo roteiro, ou muito perto dele. A segunda é que a velocidade depende do número de comandos, e não do ritmo de cada comando. A mediana do intervalo entre comandos permaneceu entre cerca de 2,5 e 3,5~s em todas as iterações, limitada pelo tempo de permanência, e os ganhos vieram de errar menos, navegar por caminhos mais curtos e aproveitar as sugestões. A terceira é que o preenchimento automático é o recurso que mais aumenta a velocidade de comunicação. No T2, em que todas as palavras foram inseridas pela barra de sugestões, a economia de comandos chegou a 79,4\%, próxima dos 81,4\% esperados, e a velocidade chegou a cerca de 3,0 palavras por minuto, contra no máximo 0,69 na digitação letra a letra do T1. Esse valor está na mesma faixa de outros sistemas de baixo custo baseados em câmera comum (Tabela~\ref{tab:comparacao_trabalhos}).

O T3, mais próximo do uso real por exigir letras digitadas e a correção de um erro antes que algumas palavras fossem sugeridas, mostrou também onde estão as principais perdas: as reversões, em que um comando é seguido do seu oposto, e as entradas na barra de sugestões sem aceitar nenhuma. Esses comandos desperdiçados explicam a diferença entre a economia obtida, de 29,9\% a 48,6\%, e a esperada, de 67,3\%, e apontam para ajustes na forma como o olhar é convertido em comandos, mais do que no modelo de linguagem.

Esses resultados devem ser lidos com as limitações discutidas na Seção~\ref{sec:res_limitacoes}: os testes tiveram um único participante, um dos autores, sem deficiência motora, e usaram três frases, duas delas do \textit{corpus} de treino. Eles mostram que o sistema funciona e que suas métricas são confiáveis, mas não substituem a avaliação com pessoas com \gls{ELA}. Ainda assim, o trabalho indica que um \textit{tablet} comum, sem nenhum equipamento adicional, pode servir de base para uma ferramenta de comunicação pelo olhar acessível a quem não tem acesso aos dispositivos comerciais.

\section{Trabalhos Futuros}

A partir dos resultados e das limitações deste trabalho, três direções de continuidade se destacam:

\begin{itemize}
    \item \textbf{Testes com usuários reais:} o passo mais importante é avaliar o aplicativo com o público-alvo. Testes com pessoas com \gls{ELA}, com mais participantes, frases fora do \textit{corpus} de treino e períodos de uso livre, permitiriam medir a velocidade e o esforço reais, observar o efeito do cansaço e da postura sobre o rastreamento e verificar se as sugestões correspondem ao que os usuários de fato querem dizer. O registro de métricas já permite esse tipo de análise, inclusive no uso livre, sem frase-alvo.
    \item \textbf{Aviso sonoro declarando a tecla selecionada:} durante os testes, foi difícil saber qual tecla estava em foco, porque, para selecionar, o olhar precisa ficar no centro da tela, longe do teclado e da mensagem. A síntese de voz, que já funciona no próprio aparelho, poderia anunciar a tecla em foco ou a letra recém-escolhida, permitindo conferir a seleção sem desviar o olhar. O aviso precisa ser curto, para não atrasar o comando seguinte.
    \item \textbf{Sistema de autocompletar com base em contextos de uso:} o modelo atual sugere palavras a partir de um único \textit{corpus}, igual para todas as situações. As sugestões poderiam considerar o contexto em que o aplicativo é usado, como o horário, o ambiente ou o interlocutor, e aprender com as frases efetivamente faladas por cada pessoa, armazenadas apenas no aparelho, de modo a priorizar o vocabulário mais provável em cada momento.
\end{itemize}
